\section{Bit complexity and the verification boundary}\label{sec:verification}

\subsection{Binary cocycle values stay controlled}

\begin{lemma}[Height growth under a bounded number of expansions]\label{lem:bits}
Let $D_0$ be the maximum absolute source edge value. After any trace with
$u$ upward events, every signed edge value has absolute value at most
$2^uD_0$. A normalized tetrahedral height row and every local coherent
normal coordinate have bit length at most $B+u+O(1)$ when $D_0<2^B$.
\end{lemma}
\begin{proof}
A downward replacement removes its central edge and creates no new global
edge. Surviving edge values are unchanged. For an upward replacement, the
new diagonal joins the formal apices $a,b$. Choose any belt vertex $v$.
Its value is
\[
 h_b-h_a=(h_b-h_v)+(h_v-h_a),
\]
the sum of two old signed edge values. Its magnitude is at most twice the
previous maximum. All other values persist. Induction proves the bound.
Each normalized local height is an edge difference from the first corner.
The coherent normal coordinates are gaps between ordered local heights,
and each gap is bounded by the local span, itself an edge difference.
\end{proof}

This bound charges arithmetic to the number of expansions, not to the
potentially enormous number of normal discs. Multiplying the initial heights
by a large integer changes the arithmetic bit size without requiring expanded
normal sheets. The tests include a 2,048-bit marking control.

\subsection{Move and identity costs}

There are at most $6(t+U)$ local edge occurrences and $2(t+U)$ paired faces.
Grouping edge occurrences and checking the fixed formal patterns enumerates
the possible events using polynomial work. Each accepted transition uses
the existing replacement producer and its independent relative-boundary
checker, followed by independent cochain transport replay. The implementations
may rebuild global incidence data; they are not claimed to perform every
move in constant time.

An output incarnation has a structural name consisting of its event's
consumed identities and local ports, its event kind, and its output slot.
Repeatedly printing a fully expanded ancestry expression could be very large.
The code instead interns these expressions as a directed acyclic graph with
exact tuple equality. A path creates only $O(t+r+U)$ name nodes. Integer IDs
are handles into the shared table; a cryptographic digest is not used as a
proof that two names are equal. For the mathematical encoding argument one
may fix any total order on the exact names. Within an execution, existing
interned IDs never change.

If every intermediate triangulation and height list is stored in a positive
certificate, there are at most $r+2U$ moves and each list has at most $t+U$
rows. A conservative bit-size bound is therefore
\begin{equation}\label{eq:trace-proof-size}
 O\!\left((r+2U+1)(t+U)
     (B+U+\log(t+U+1))\right),
\end{equation}
plus the terminal component certificate. The geometric-trace portion is
polynomial in the parameters of one trace; the terminal component
certificate is an additional term. The artifact uses explicit snapshots for review;
it makes no unsupported claim that all certificates of an exhaustive search
fit in one small proof. Storing every endpoint is optional.

\subsection{Independent endpoint replay}

The new endpoint certificate records the total upward allowance, the allowed
initial region, the ordered geometric transition certificates, and the final
normal coordinates. The verifier starts from the independently supplied
source triangulation and source heights. It replays each formal replacement
and cocycle transport using the existing checkers, tracks which original
tetrahedra survive, counts upward events, and reconstructs the final local
coordinates. It does not invoke commitment search, sleep sets, event generation,
the cocycle seed solver, or a move producer.

The initial-region claim is checked by tracking original incarnations through
the canonical survivor order in every recorded transition. A trace that
consumes a protected original tetrahedron is rejected even if its final
triangulation and surface would otherwise be valid. A false upward budget
is likewise rejected. The search's optimization bookkeeping is not trusted
as a substitute for these checks.

\subsection{Disc components and diagram provenance}

For a supplied endpoint surface, the existing compressed normal-component
kernel determines whether it contains an essential compressing disc, retaining
a replayable certificate. Binary-coordinate component computations build on
the interval-orbit approach of Agol--Hass--Thurston \cite{AHT} and the project's
later checked implementation. The local search theorem can use any polynomial
endpoint predicate; if a selected implementation has a different proved
complexity, that cost must be charged separately.

No inference is made from total Euler characteristic alone. A disconnected
surface can contain spheres that raise its total Euler characteristic while
leaving its other components unsuitable. Likewise, preserving a cohomology
class does not mean preserving the topology of its coherent fibre. The
component predicate is the relevant final test.

The new \code{pachner\_seed\_decide} adapter authenticates an original knot
diagram using the maintained diagram-to-exterior certificate. It optionally
replays boundary shellings, constructs the starting integral cocycle, runs
the bounded search, and converts a successful endpoint into the unchanged
\code{diagram-transport-disc-v1} format. The existing independent
\code{verify\_transport\_disk\_certificate} checks the entire chain. A
verified compressing disc in the authenticated knot exterior establishes
the positive unknot conclusion through the established geometric criterion.

The statuses have the following meanings:
\begin{center}
\small
\begin{tabularx}{\textwidth}{@{}lX@{}}
\toprule
Status & Exact meaning\\
\midrule
\status{DISC\_FOUND} & A source-bound endpoint contains a verified compressing-disc component\\
\status{ENDPOINT\_SELECTED} & The caller selected an endpoint; no extra topological claim is implied\\
\status{COMPLETE\_BOUNDED\_FAMILY} & The declared single-region or unrestricted finite family was exhausted\\
\status{COMPLETE\_BOUNDED\_REGIONS} & Every permitted initial region and its finite family was exhausted\\
\status{INCONCLUSIVE} & A required search or component query was capped, or the diagram adapter found no verified disc\\
\status{UNKNOT} & The diagram adapter's complete source-to-disc certificate independently replayed\\
\bottomrule
\end{tabularx}
\end{center}

The complete-family statuses are search results, not independently supplied
negative unknot certificates. In particular, the diagram adapter maps bounded
family exhaustion to \status{INCONCLUSIVE}. A caller can then invoke the
maintained complete recognizer under its own remaining budget.

\subsection{Resource limits and observational guarantees}

The node allowance counts popped search frames. Commitment assignment frames
and geometric trace frames are different quantities, so their node caps are
not equal units of geometric coverage. The work allowance is a cooperative
guard counter shared by preparation, moves, queries and positive replay;
it is not a bit-operation counter. Disc-query caps apply at individual
endpoints, and any incomplete required disc query prevents a final claim
of exhausted disc search.

Callbacks receive private certificate copies. Mutation cannot alter the
search state. Caller exceptions are shielded through internal cap and
malformed-certificate handlers and are re-raised with their original
identity. These details matter for faithful integration into an existing
recognition budget.
